% ---------------------------------------------------------------- 7 limitations, conclusion
\section{Limitations}\label{sec:limits}
Both backbones make confident errors on the blind set (\BlindWrong{} and \UPBlindWrong{}), almost
all on drawings with nearly invisible gaps: confidence is calibrated against view-dependent
ambiguity, not against consistent misreading. The picture's own gap evidence adds answers on a
second blind set without adding errors, but it does not remove those errors. The Minkowski
variants rest on three paired seeds per device, and the plain model was not trained on the GPU,
so their effect size is uncertain. All quantitative cable results are on renders. No annotated photographs of real cables
exist yet. On public pictures the system answers drawings and renders but no photograph
of a physical object, and those pictures informed our choice of training styles. Strands
that touch (gap below one \cw) and double crossings smaller than a few pixels
cannot be resolved from one view; multi-view fusion is implemented but untested on real
data. The output is a knot type, not a physical untangling plan: friction, stiffness
and tension decide how hard a cable is to undo.

\section{Conclusion}\label{sec:concl}
Reading the diagram and computing the topology, instead of classifying the image,
gives a knot recogniser that generalises past a census, explains its answers and
abstains when the image is not enough. The gain comes from the formulation, not from a
particular backbone: standard U-Nets do better than our first transformer, and a U-Net with a
transformer bottleneck does better than both on blind and real pictures. On renders it reads crossings and \ou{} almost
perfectly and makes no wrong call. The open problem is perception on real photographs,
for which we release the benchmark protocol, annotation tool and scoring code.
